\documentclass[12pt,a4paper]{article}
\usepackage[UTF8]{ctex}
\usepackage{amsmath,amssymb,amsthm}
\usepackage{geometry}

\geometry{margin=2.5cm}

\title{任务空间动力学方程中"执行器提供零力和力矩"的含义}
\author{第8章补充说明}
\date{}

\begin{document}

\maketitle

\section{问题提出}

\textbf{原文}：
\begin{quote}
These are the dynamic equations expressed in end-effector frame coordinates. If an external wrench $F$ is applied to the end-effector frame then, assuming the actuators provide zero forces and torques, the motion of the end-effector frame is governed by these equations.
\end{quote}

\textbf{中文翻译}：
\begin{quote}
这些是用末端执行器坐标系表示的动力学方程。如果对末端执行器坐标系施加外部wrench $F$，那么，假设执行器提供零力和力矩，末端执行器坐标系的运动由这些方程决定。
\end{quote}

\textbf{问题}：如何理解"假设执行器提供零力和力矩"？

\section{任务空间动力学方程回顾}

\subsection{方程形式}

\textbf{任务空间动力学方程}：
\begin{equation}
F = \Lambda(\theta)\dot{V} + \eta(\theta, V) \tag{8.89}
\end{equation}

其中：
\begin{itemize}
\item $F$：在末端执行器坐标系中表示的wrench（6×1向量，包含力和力矩）
\item $\Lambda(\theta)$：任务空间质量矩阵（6×6）
\item $\dot{V}$：末端执行器加速度（twist的导数，6×1）
\item $\eta(\theta, V)$：任务空间的科里奥利力、向心力和重力项（6×1）
\end{itemize}

\subsection{方程的物理意义}

\textbf{这个方程描述了}：
\begin{itemize}
\item 要产生加速度$\dot{V}$，需要施加wrench $F$
\item $F$包括两部分：
\begin{enumerate}
\item 惯性力：$\Lambda(\theta)\dot{V}$（产生加速度所需的力）
\item 科里奥利、向心力和重力：$\eta(\theta, V)$（由于速度和位置产生的力）
\end{enumerate}
\end{itemize}

\section{"执行器提供零力和力矩"的含义}

\subsection{关节空间动力学}

\textbf{关节空间动力学方程}：
\begin{equation}
\tau = M(\theta)\ddot{\theta} + h(\theta, \dot{\theta}) \tag{8.82}
\end{equation}

其中：
\begin{itemize}
\item $\tau$：关节力矩（由执行器提供）
\item $M(\theta)$：关节空间质量矩阵
\item $\ddot{\theta}$：关节加速度
\item $h(\theta, \dot{\theta})$：科里奥利、向心力和重力项
\end{itemize}

\subsection{执行器提供零力矩的情况}

\textbf{假设}：$\tau = 0$（执行器不提供任何力矩）

\textbf{这意味着}：
\begin{itemize}
\item 所有关节的执行器都不产生力矩
\item 机器人处于"被动"状态
\item 只有外力作用在机器人上
\end{itemize}

\textbf{关节空间动力学方程变为}：
\begin{equation}
0 = M(\theta)\ddot{\theta} + h(\theta, \dot{\theta}) - J^T(\theta)F_{\text{ext}}
\end{equation}

其中$F_{\text{ext}}$是作用在末端执行器上的外部wrench。

\textbf{重新整理}：
\begin{equation}
J^T(\theta)F_{\text{ext}} = M(\theta)\ddot{\theta} + h(\theta, \dot{\theta}) \tag{*}
\end{equation}

\subsection{转换到任务空间}

\textbf{从关节空间到任务空间}：

使用关系：
\begin{align}
V &= J(\theta)\dot{\theta} \tag{8.83} \\
\dot{V} &= \dot{J}(\theta)\dot{\theta} + J(\theta)\ddot{\theta} \tag{8.84}
\end{align}

\textbf{推导}：

从方程(*)开始：
\begin{equation}
J^T(\theta)F_{\text{ext}} = M(\theta)\ddot{\theta} + h(\theta, \dot{\theta})
\end{equation}

左乘$J^{-T}(\theta)$：
\begin{equation}
F_{\text{ext}} = J^{-T}(\theta)M(\theta)\ddot{\theta} + J^{-T}(\theta)h(\theta, \dot{\theta})
\end{equation}

从$\dot{V} = \dot{J}(\theta)\dot{\theta} + J(\theta)\ddot{\theta}$，得到：
\begin{equation}
\ddot{\theta} = J^{-1}(\theta)(\dot{V} - \dot{J}(\theta)\dot{\theta}) = J^{-1}(\theta)(\dot{V} - \dot{J}(\theta)J^{-1}(\theta)V)
\end{equation}

代入并整理，得到：
\begin{equation}
F_{\text{ext}} = \Lambda(\theta)\dot{V} + \eta(\theta, V) \tag{8.89}
\end{equation}

其中：
\begin{align}
\Lambda(\theta) &= J^{-T}(\theta)M(\theta)J^{-1}(\theta) \tag{8.90} \\
\eta(\theta, V) &= J^{-T}(\theta)h(\theta, J^{-1}(\theta)V) - \Lambda(\theta)\dot{J}(\theta)J^{-1}(\theta)V \tag{8.91}
\end{align}

\section{物理意义的理解}

\subsection{关键点}

\textbf{当$\tau = 0$时}：
\begin{itemize}
\item 执行器不提供任何力矩
\item 机器人完全由外力$F_{\text{ext}}$驱动
\item 末端执行器的运动完全由方程(8.89)决定
\end{itemize}

\textbf{方程(8.89)的含义}：
\begin{itemize}
\item 如果对末端执行器施加外力$F_{\text{ext}}$
\item 且执行器不提供力矩（$\tau = 0$）
\item 那么末端执行器会产生加速度$\dot{V}$
\item 这个加速度由方程(8.89)决定
\end{itemize}

\subsection{具体例子}

\textbf{例子1：自由落体}

\textbf{场景}：
\begin{itemize}
\item 机器人悬挂在空中
\item 执行器不提供力矩（$\tau = 0$）
\item 只有重力作用在末端执行器上
\end{itemize}

\textbf{外力}：
\begin{equation}
F_{\text{ext}} = \begin{bmatrix} 0 \\ 0 \\ -mg \\ 0 \\ 0 \\ 0 \end{bmatrix}
\end{equation}

（假设重力沿$z$轴负方向）

\textbf{运动方程}：
\begin{equation}
F_{\text{ext}} = \Lambda(\theta)\dot{V} + \eta(\theta, V)
\end{equation}

\textbf{结果}：
\begin{itemize}
\item 末端执行器会加速下落
\item 加速度$\dot{V}$由方程决定
\item 这类似于自由落体，但考虑了机器人的动力学
\end{itemize}

\textbf{例子2：外力推动}

\textbf{场景}：
\begin{itemize}
\item 机器人静止
\item 执行器不提供力矩（$\tau = 0$）
\item 用户用手推动末端执行器，施加力$f_{\text{hand}}$
\end{itemize}

\textbf{外力}：
\begin{equation}
F_{\text{ext}} = \begin{bmatrix} f_{\text{hand},x} \\ f_{\text{hand},y} \\ f_{\text{hand},z} \\ 0 \\ 0 \\ 0 \end{bmatrix}
\end{equation}

\textbf{运动方程}：
\begin{equation}
F_{\text{ext}} = \Lambda(\theta)\dot{V} + \eta(\theta, V)
\end{equation}

\textbf{结果}：
\begin{itemize}
\item 末端执行器会加速运动
\item 加速度$\dot{V}$由方程决定
\item 这类似于推动一个质量-弹簧-阻尼器系统
\end{itemize}

\section{为什么这个假设很重要？}

\subsection{理解动力学方程}

\textbf{这个假设帮助我们理解}：
\begin{itemize}
\item 动力学方程描述的是"如果施加力$F$，会产生什么运动"
\item 当$\tau = 0$时，只有外力$F_{\text{ext}}$作用
\item 运动完全由动力学方程决定
\item 这类似于牛顿第二定律：$F = ma$
\end{itemize}

\subsection{在控制中的应用}

\textbf{在阻抗控制中}：
\begin{itemize}
\item 我们希望机器人表现得像质量-弹簧-阻尼器系统
\item 方程(8.89)描述了机器人的"自然"动力学
\item 通过控制律，我们可以修改这个动力学
\item 但基础仍然是方程(8.89)
\end{itemize}

\textbf{在力控制中}：
\begin{itemize}
\item 我们希望控制末端执行器施加的力
\item 方程(8.89)告诉我们，要产生某个加速度，需要多大的力
\item 通过控制律，我们可以产生期望的力
\end{itemize}

\section{与有执行器力矩的情况对比}

\subsection{有执行器力矩的情况}

\textbf{关节空间动力学}：
\begin{equation}
\tau = M(\theta)\ddot{\theta} + h(\theta, \dot{\theta}) - J^T(\theta)F_{\text{ext}} \tag{8.82-full}
\end{equation}

\textbf{转换到任务空间}：
\begin{equation}
J^T(\theta)\tau + F_{\text{ext}} = \Lambda(\theta)\dot{V} + \eta(\theta, V)
\end{equation}

\textbf{重新整理}：
\begin{equation}
F_{\text{total}} = \Lambda(\theta)\dot{V} + \eta(\theta, V)
\end{equation}

其中$F_{\text{total}} = J^T(\theta)\tau + F_{\text{ext}}$是总wrench（执行器产生的wrench + 外部wrench）。

\subsection{对比}

\textbf{当$\tau = 0$时}：
\begin{equation}
F_{\text{ext}} = \Lambda(\theta)\dot{V} + \eta(\theta, V)
\end{equation}

\textbf{当$\tau \neq 0$时}：
\begin{equation}
J^T(\theta)\tau + F_{\text{ext}} = \Lambda(\theta)\dot{V} + \eta(\theta, V)
\end{equation}

\textbf{区别}：
\begin{itemize}
\item 当$\tau = 0$时，只有外力$F_{\text{ext}}$驱动运动
\item 当$\tau \neq 0$时，执行器力矩$J^T(\theta)\tau$也驱动运动
\item 总wrench $F_{\text{total}} = J^T(\theta)\tau + F_{\text{ext}}$决定运动
\end{itemize}

\section{数学推导的完整性}

\subsection{从关节空间到任务空间}

\textbf{完整推导}：

\textbf{步骤1：关节空间动力学}
\begin{equation}
\tau = M(\theta)\ddot{\theta} + h(\theta, \dot{\theta}) - J^T(\theta)F_{\text{ext}}
\end{equation}

\textbf{步骤2：假设$\tau = 0$}
\begin{equation}
0 = M(\theta)\ddot{\theta} + h(\theta, \dot{\theta}) - J^T(\theta)F_{\text{ext}}
\end{equation}

\textbf{步骤3：重新整理}
\begin{equation}
J^T(\theta)F_{\text{ext}} = M(\theta)\ddot{\theta} + h(\theta, \dot{\theta})
\end{equation}

\textbf{步骤4：左乘$J^{-T}(\theta)$}
\begin{equation}
F_{\text{ext}} = J^{-T}(\theta)M(\theta)\ddot{\theta} + J^{-T}(\theta)h(\theta, \dot{\theta})
\end{equation}

\textbf{步骤5：代入$\ddot{\theta} = J^{-1}(\theta)(\dot{V} - \dot{J}(\theta)J^{-1}(\theta)V)$}
\begin{equation}
F_{\text{ext}} = J^{-T}(\theta)M(\theta)J^{-1}(\theta)(\dot{V} - \dot{J}(\theta)J^{-1}(\theta)V) + J^{-T}(\theta)h(\theta, J^{-1}(\theta)V)
\end{equation}

\textbf{步骤6：展开并整理}
\begin{align}
F_{\text{ext}} &= J^{-T}(\theta)M(\theta)J^{-1}(\theta)\dot{V} \\
&\quad - J^{-T}(\theta)M(\theta)J^{-1}(\theta)\dot{J}(\theta)J^{-1}(\theta)V \\
&\quad + J^{-T}(\theta)h(\theta, J^{-1}(\theta)V)
\end{align}

\textbf{步骤7：定义$\Lambda(\theta)$和$\eta(\theta, V)$}
\begin{align}
\Lambda(\theta) &= J^{-T}(\theta)M(\theta)J^{-1}(\theta) \\
\eta(\theta, V) &= J^{-T}(\theta)h(\theta, J^{-1}(\theta)V) - \Lambda(\theta)\dot{J}(\theta)J^{-1}(\theta)V
\end{align}

\textbf{步骤8：得到最终方程}
\begin{equation}
F_{\text{ext}} = \Lambda(\theta)\dot{V} + \eta(\theta, V) \tag{8.89}
\end{equation}

\section{总结}

\subsection{关键要点}

\begin{enumerate}
\item \textbf{任务空间动力学方程}：
\begin{equation}
F = \Lambda(\theta)\dot{V} + \eta(\theta, V)
\end{equation}
描述要产生加速度$\dot{V}$，需要施加wrench $F$。

\item \textbf{"执行器提供零力和力矩"的含义}：
\begin{itemize}
\item 假设$\tau = 0$（所有执行器不产生力矩）
\item 只有外力$F_{\text{ext}}$作用在末端执行器上
\item 末端执行器的运动完全由方程(8.89)决定
\end{itemize}

\item \textbf{物理意义}：
\begin{itemize}
\item 这类似于牛顿第二定律：$F = ma$
\item 如果施加力$F$，会产生加速度$\dot{V}$
\item 但这里考虑了机器人的复杂动力学（质量矩阵、科里奥利力等）
\end{itemize}

\item \textbf{应用}：
\begin{itemize}
\item 理解机器人的"自然"动力学
\item 在阻抗控制中，修改这个动力学
\item 在力控制中，计算需要的力
\end{itemize}
\end{enumerate}

\subsection{公式总结}

\textbf{任务空间动力学方程}（当$\tau = 0$时）：
\begin{equation}
F_{\text{ext}} = \Lambda(\theta)\dot{V} + \eta(\theta, V)
\end{equation}

\textbf{其中}：
\begin{align}
\Lambda(\theta) &= J^{-T}(\theta)M(\theta)J^{-1}(\theta) \\
\eta(\theta, V) &= J^{-T}(\theta)h(\theta, J^{-1}(\theta)V) - \Lambda(\theta)\dot{J}(\theta)J^{-1}(\theta)V
\end{align}

\textbf{物理意义}：
\begin{itemize}
\item 如果对末端执行器施加外力$F_{\text{ext}}$
\item 且执行器不提供力矩（$\tau = 0$）
\item 那么末端执行器会产生加速度$\dot{V}$
\item 这个加速度由方程决定
\end{itemize}

\end{document}

